\chapter{회로도 없이 기계를 고치는 법}
% full version: chapters/17_debugging.tex

\pencilfig{17}{\pencilart{17}}{회로도 없는 기판. 프로브는 수십억 개의 전선 가운데 몇 개만 듣는다.}

동작하는 기판을 회로도 없이 받아서 고쳐 달라는 부탁을 받았다고 하자. 엔지니어에게는 도구가 네 가지 있다. 신호를 들어 보는 프로브, 자기 신호를 넣어 주는 신호 발생기, 부품 하나를 떼어 내고 무엇이 망가지는지 보는 방법, 그리고 제어 레지스터에 대한 접근이다. 뇌도 바로 이 도구들로 연구한다.

\section*{뉴런 천 개를 듣는 프로브}

뇌에 꽂은 전극은 가까이 있는 뉴런 몇 개의 딸깍을 듣는다. 860억 개 가운데 전극 천 개라면 아무것도 아닌 것 같다. 그런데 이것으로 화면의 커서를 움직이기에는 충분하다. 왜일까? 이웃한 뉴런들은 비슷하게 요동하고, 팔을 움직이는 활동은 실제로 공통 변수 열두어 개로 기술된다. 모두를 들을 필요는 없다. 이 몇 안 되는 “목소리”만 들으면 된다.

전극 천 개에서 나오는 원시 신호는 초당 수억 비트다. 두개골 아래에서 무선 채널로 내보내기에는 너무 많다. 반면 딸깍 자체가 나르는 정보는 그보다 천 배쯤 적다. 그래서 딸깍을 임플란트 안에서 바로 골라내고 그것만 밖으로 보내는 편이 이득이다. 눈에서 본 것과 같은 압축 기법이다.

\section*{뉴럴링크가 하는 일}

뉴럴링크(Neuralink)는 이 문제의 공학적인 측면에 뛰어들었다. 단단한 바늘 대신 가늘고 유연한 실, 실을 심는 로봇, 밀봉된 무선 임플란트다. 이 회사가 발표한 논문에는 케이블로 컴퓨터에 연결된 시스템이 기술되어 있고, 압축과 무선 통신은 계획으로 언급되어 있다. 팔이 마비된 사람이 커서를 조종한다는 데이터는 아직 주로 회사 자체의 보고로 알려져 있다. 그 보고에 따르면 임플란트에는 실 수십 가닥에 걸쳐 채널이 천 개쯤 있고, 심은 뒤 실 일부가 제자리에서 밀려났다. 발상 자체는 새롭지 않다. 전극 100개 정도의 단단한 전극 배열을 쓰는 학계 연구진은 이미 사람들이 커서를 조종하고, 생각만으로 글자를 “쓰고”, 말하려는 시도를 분당 60단어 정도의 속도로 화면의 글로 바꾸게 해 주었다.

\section*{읽기와 쓰기}

디코더는 많은 뉴런의 딸깍 빈도를 읽어 초당 10비트도 안 되는 정보를 얻는다. 뉴런 자체가 나르는 것에 비하면 아주 작은 몫이다. 그 대신 뇌와 디코더는 서로에게서 배운다. 사람은 디코더에 적응하고, 디코더는 사람에게 적응한다. 서로에게 맞춰 가는 학습자 둘은 까다로운 문제다. 각자는 안정해도 함께 있으면 흔들릴 수 있다.

뇌에 쓰는 것은 읽는 것보다 어렵다. 전극을 지나는 전류는 원하는 뉴런 하나가 아니라 주변 전체를 흥분시킨다. 빛나는 픽셀처럼 빛의 점을 “그릴” 수는 있고, 그러면 사람은 간단한 글자를 알아본다. 그러나 뇌가 그림을 보도록 눈의 압축 부호를 써 넣는 방법은 아직 아무도 모른다.

\section*{프로브가 보지 못하는 것}

전극은 데이터가 흐르는 전화망을 듣는다. 그러나 방송국, 곧 도파민, 세로토닌, 노르아드레날린, 아세틸콜린은 물질의 농도이며, 프로브는 이것을 듣지 못한다. 기억이 모두 저장된 연결의 세기도 보지 못한다. 사람이 느끼는 그대로의 통증도 보지 못한다. 데이터는 보지만 레지스터는 보지 못하는 디버거는 같은 입력이 왜 다른 응답을 내는지 언제까지나 의아해할 것이다.

\fullversion{17}
